\documentclass{article}
\usepackage[utf8]{inputenc}
\usepackage[spanish]{babel}
\usepackage{tikz}
\usepackage{geometry}
\usepackage{graphicx} % Requerido para \resizebox
\geometry{a4paper, margin=1in}

\usetikzlibrary{shapes.geometric, arrows.meta, positioning, calc}

\title{Código TikZ para la Arquitectura de MTDE-Net (Tesis)}
\author{Proyecto FLIR Image Processing}
\date{\today}

\begin{document}

\maketitle

Este documento contiene el código de diagramación vectorial en TikZ listo para compilar en Overleaf o cualquier distribución de \LaTeX.

Para evitar que los diagramas excedan el ancho de la página, se ha envuelto cada entorno \texttt{tikzpicture} dentro de un comando \texttt{\textbackslash resizebox\{\textbackslash textwidth\}\{!\}\{...\}}. Esto escalará automáticamente el diagrama para ajustarse perfectamente a los márgenes de tu documento.

---

\section{Diagrama 1: Arquitectura General de MTDE-Net}

Copia este bloque de código en tu tesis para graficar el modelo general:

\begin{figure}[htbp]
\centering
% \resizebox fuerza al diagrama a ajustarse al ancho del texto de la hoja
\resizebox{\textwidth}{!}{
\begin{tikzpicture}[
    auto,
    >=Stealth,
    node distance=1.2cm,
    % Definición de estilos de cajas
    block/.style={draw, fill=blue!8, rectangle, minimum height=2.8em, minimum width=6.2em, rounded corners, align=center, text width=6em},
    input/.style={draw, fill=red!8, ellipse, align=center, text width=6em, minimum height=2.8em},
    output/.style={draw, fill=green!8, rounded corners, align=center, text width=6em, minimum height=2.8em},
    fusion/.style={draw, fill=yellow!15, circle, minimum size=1.3cm, align=center, inner sep=0pt}
]
    % --- Rama Visual (Arriba) ---
    \node [input] (img) {Imagen\\ Térmica\\ \small(1x112x112)};
    \node [block, right=of img] (stem) {Stem Conv\\ \small(32x56x56)};
    \node [block, right=of stem] (stages) {RAB Stages (1-4)\\ \small(256x7x7)};
    \node [block, right=of stages] (gap) {Global Avg\\ Pool (GAP)\\ \small(256)};

    % --- Rama Tabular (Abajo) ---
    \node [input, below=of img, yshift=-1.0cm] (tab) {Metadatos\\ Tabulares\\ \small(10 variables)};
    \node [block, right=of tab] (mlp) {MLP Tabular\\ \small(64)};

    % --- Fusión y Cabezal ---
    \node [fusion, right=of gap, xshift=0.8cm] (concat) {Late\\ Fusion\\ \small(320)};
    \node [block, right=of concat] (head) {Cabezal\\ Regresión\\ \small(256)};
    \node [output, right=of head] (out) {Predicción\\ \small Tiempo (s)};

    % --- Conexiones y Flechas ---
    \draw [->, thick] (img) -- (stem);
    \draw [->, thick] (stem) -- (stages);
    \draw [->, thick] (stages) -- (gap);
    \draw [->, thick] (tab) -- (mlp);
    
    % Conectar GAP a la Fusión
    \draw [->, thick] (gap) -- (concat);
    
    % Conectar MLP a la Fusión con codo ortogonal (en forma de L)
    \draw [->, thick] (mlp.east) -| ($(concat.west)-(0.5,0)$) |- (concat.south west);

    \draw [->, thick] (concat) -- (head);
    \draw [->, thick] (head) -- (out);
\end{tikzpicture}
}
\caption{Diagrama de bloques de la arquitectura multimodal MTDE-Net con fusión tardía.}
\label{fig:mtde_net_architecture}
\end{figure}

---

\newpage
\section{Diagrama 2: Estructura Interna del Residual Attention Block (RAB)}

Copia este bloque de código para explicar la atención dual dentro de cada bloque residual:

\begin{figure}[htbp]
\centering
% \resizebox para ajustar el bloque al ancho de la hoja
\resizebox{\textwidth}{!}{
\begin{tikzpicture}[
    auto,
    >=Stealth,
    node distance=1.5cm,
    % Estilos para bloques internos del RAB
    block/.style={draw, fill=blue!8, rectangle, minimum height=2.8em, minimum width=6.5em, rounded corners, align=center, text width=6.2em},
    sum/.style={draw, fill=yellow!15, circle, minimum size=0.6cm, align=center, inner sep=0pt}
]
    % --- Camino Principal ---
    \node (in) {Entrada $C_{in}$};
    \node [block, right=of in] (conv1) {Conv 3x3\\ + BN + Act};
    \node [block, right=of conv1] (conv2) {Conv 3x3\\ + BN};
    
    % --- Ramas de Atención (SA + CA) ---
    \node [block, above right=of conv2, yshift=-0.4cm, xshift=0.3cm] (sa) {Atención\\ Espacial (SA)};
    \node [block, below right=of conv2, yshift=0.4cm, xshift=0.3cm] (ca) {Atención\\ Canal (CA)};
    
    \node [sum, right=of conv2, xshift=3.5cm] (sum_att) {+};
    \node [sum, right=of sum_att] (sum_skip) {+};
    
    \node [block, right=of sum_skip] (act) {PReLU\\ Final};
    \node [right=of act] (out) {Salida $C_{out}$};

    % --- Conexiones de Flujo ---
    \draw [->, thick] (in) -- (conv1);
    \draw [->, thick] (conv1) -- (conv2);
    
    % Bifurcación hacia la atención dual (SA y CA)
    \draw [->, thick] (conv2.east) -- +(0.4,0) |- (sa.west);
    \draw [->, thick] (conv2.east) -- +(0.4,0) |- (ca.west);
    
    % Sumar los mapas de atención
    \draw [->, thick] (sa.east) -| (sum_att.north);
    \draw [->, thick] (ca.east) -| (sum_att.south);
    
    \draw [->, thick] (sum_att) -- (sum_skip);
    
    % --- Conexión de Salto (Skip Connection) ---
    % Dibuja la flecha por abajo para evitar cruzar otras líneas
    \draw [->, thick] (in.east) -- +(0.3,0) -- +(0.3,-2.2) -| (sum_skip.south);
    
    \draw [->, thick] (sum_skip) -- (act);
    \draw [->, thick] (act) -- (out);
\end{tikzpicture}
}
\caption{Estructura interna del Residual Attention Block (RAB) con ramas duales de atención.}
\label{fig:rab_block}
\end{figure}

\end{document}
